\BourbakiExercises{习题}

\begin{enumerate}
\def\labelenumi{\arabic{enumi})}
\tightlist
\item
  设 M 是一个 A-模，N 是 M 的一个直因子子模。
\end{enumerate}

\begin{enumerate}
\def\labelenumi{\alph{enumi})}
\item
  证明 N 在 M 的双交换子 \(A_{M}^{\prime\prime}\) 下稳定。
\item
  证明每个从 N 到 M 的 A-线性映射都是 \(A_{M}^{\prime\prime}\)-线性的。
\item
  证明在 N 上的限制定义一个环同态 \(A_{M}^{\prime\prime} \to A_{N}^{\prime\prime}\)。
\item
  假设 M/N 是一些同构于 N 的商的子模之和。证明典范同态 \(A_{M}^{\prime\prime} \rightarrow A_{N}^{\prime\prime}\) 是单的。若 N 是平衡的，则 M 亦然。
\end{enumerate}

\begin{enumerate}
\def\labelenumi{\arabic{enumi})}
\setcounter{enumi}{1}
\tightlist
\item
  设 K 是一个交换域（域），A 是 \(\mathbf{M}_{3}(\mathrm{K})\) 的由形如下式的矩阵组成的子环
\end{enumerate}

\[
\left( \begin{array}{c c c} a & 0 & 0 \\ b & c & 0 \\ 0 & 0 & a \end{array} \right)
\]

其中 \(a, b, c\) 属于 K。设 M 是 A-模 \(\mathrm{K}^3\)；它是子模 \(\mathrm{N} = \mathrm{Ke}_1 + \mathrm{Ke}_2\) 与 \(\mathrm{P} = \mathrm{Ke}_3\) 的直和。

\begin{enumerate}
\def\labelenumi{\alph{enumi})}
\item
  证明 M 是平衡的，典范映射 \(A_{M}^{\prime\prime} \rightarrow A_{N}^{\prime\prime}\)（习题 1）是单的但不是满的，而典范映射 \(A_{M}^{\prime\prime} \rightarrow A_{P}^{\prime\prime}\) 是满的但不是单的。
\item
  由此给出环 B 上的模 Q 及其直因子子模 R 的一个例子，使典范映射 \(B_{Q}^{\prime\prime} \rightarrow B_{R}^{\prime\prime}\) 既非单又非满（取 \(B = A \times A\) 及 \(Q = M \times M\)）。
\end{enumerate}

\begin{enumerate}
\def\labelenumi{\arabic{enumi})}
\setcounter{enumi}{2}
\item
  设 B 是一个环，I 是 B 的一个双边理想；设 A 是 \(B \times B\) 的由偶 \((x, y)\)（满足 \(x - y \in I\)）组成的子环。设 \(B_{1}\)（相应地，\(B_{2}\)）是赋以由第一个（相应地，第二个）投影导出的 A-模结构的加群 B。证明 A-模 \(B_{1}\) 与 \(B_{2}\) 都是平衡的。在 I 满足什么条件时 A-模 \(B_{1} \oplus B_{2}\) 是平衡的？
\item
  设 A 与 B 是环，P 是一个 (A, B)-双模。
\end{enumerate}

\begin{enumerate}
\def\labelenumi{\alph{enumi})}
\item
  假设 B-模 P 是生成模且 A-模 P 是平衡的。证明对每个平衡 B-模 M，A-模 \(P \otimes_{B} M\) 是平衡的。
\item
  假设 A-模 P 是生成模，且映射 \(b \mapsto b_{P}\)（从 B 到 \(\operatorname{End}_{\mathrm{A}}(\mathrm{P})\)）是双射。证明对每个生成 B-模 M，A-模 \(P \otimes_{B} M\) 是生成模。
\end{enumerate}

\begin{enumerate}
\def\labelenumi{\arabic{enumi})}
\setcounter{enumi}{4}
\item
  设 A 是一个环，M 是一个 A-模，N 是 M 的一个子模。证明若 A-模 M/N 是生成模，则 A-模 M 亦然。用一个例子证明 N 可以是生成模而 M 不是。
\item
  设 \((\mathrm{A}_i)_{i\in \mathrm{I}}\) 是一族环，A 是它们的乘积。证明 A-模 \(\bigoplus \mathrm{A}_i\) 是投射的且忠实的。它是否总是生成模？
\item
  设 A 是一个环，M 是一个 A-模。证明若 M 是生成模，则右 A-模 \(\mathrm{M}^{*} = \mathrm{Hom}_{\mathrm{A}}(\mathrm{M}, \mathrm{A}_{s})\) 是生成模。用一个例子证明 \(M^{*}\) 可以是生成模而 M 不是（参见习题 6）。
\item
  设 A 是一个主理想整环，M 是一个 A-模。证明下列性质等价：
\end{enumerate}

\begin{enumerate}
\def\labelenumi{(\roman{enumi})}
\item
  A-模 M 是生成模。
\item
  M 的对偶非零。
\item
  模 M 含有一个同构于 \(A_{s}\) 的直因子子模。
\end{enumerate}

\begin{enumerate}
\def\labelenumi{\arabic{enumi})}
\setcounter{enumi}{8}
\tightlist
\item
  设 \(D\) 是一个体，\(V\) 是 \(D\) 上一个右向量空间。我们把 \(V\) 视为环 \(A = \operatorname{End}_D(V)\) 上的一个左模。
\end{enumerate}

\begin{enumerate}
\def\labelenumi{\alph{enumi})}
\item
  证明 A-模 V 是投射的、有限生成的、忠实的且平衡的。证明 (D,A)-双模 \(\mathrm{Hom}_{\mathrm{A}}(\mathrm{V},\mathrm{A}_{s})\) 与 \(\mathrm{V}^{*}=\mathrm{Hom}_{\mathrm{D}}(\mathrm{V},\mathrm{D}_{d})\) 典范同构。
\item
  证明 A-模 V 是生成模当且仅当 D-向量空间 V 的维数有限。
\item
  A-模 V 的迹理想是什么？
\end{enumerate}

\begin{enumerate}
\def\labelenumi{\arabic{enumi})}
\setcounter{enumi}{9}
\tightlist
\item
  设 M 是一个 A-模。证明下列性质等价：
\end{enumerate}

\begin{enumerate}
\def\labelenumi{(\roman{enumi})}
\item
  A-模 M 是有限生成的。
\item
  对每个生成 A-模 P，存在一个自然数 n 及一个从 \(P^{n}\) 到 M 的满同态 f。
\end{enumerate}

\begin{enumerate}
\def\labelenumi{\arabic{enumi})}
\setcounter{enumi}{10}
\tightlist
\item
  设 \(\mathrm{P}\) 是一个有限生成投射 A-模；把它的迹理想记作 \(\tau(\mathrm{P})\)。\\
\end{enumerate}

\begin{enumerate}
\def\labelenumi{\alph{enumi})}
\item
  证明我们有 \(\tau(\mathrm{P})^2 = \tau(\mathrm{P})\) 及 \(\tau(\mathrm{P})\mathrm{P} = \mathrm{P}\)。
\item
  证明典范映射 \(\tau(\mathrm{P})\otimes_{\mathrm{A}}\mathrm{P}\to\mathrm{P}\) 是同构。
\end{enumerate}

¶ 12) 设 A 是一个环，P 是一个有限生成 A-模。证明下列性质等价：

\begin{enumerate}
\def\labelenumi{(\roman{enumi})}
\item
  A-模 P 是投射的且是生成模。
\item
  A-模 \(\mathrm{P}^{(\mathbf{N})}\) 与 \(\mathrm{A}_s^{(\mathbf{N})}\) 同构。
\end{enumerate}

（为证蕴含关系 \((\mathrm{i})\Rightarrow(\mathrm{ii})\)，用归纳法构造整数的严格递增序列 \((n_{k})\) 与 \((m_{k})\) 以及 A-模的满同态 \(f_{k}:A_{s}^{n_{k}}\to P^{m_{k}}\) 与 \(g_{k}:P^{m_{k+1}}\to A_{s}^{n_{k}}\)，使得 \(g_{k-1}\circ f_{k}\) 与 \(f_{k}\circ g_{k}\) 是典范投影。）

\begin{enumerate}
\def\labelenumi{\arabic{enumi})}
\setcounter{enumi}{12}
\item
  设 k 是一个环，n 是一个整数 \(\geqslant 2\)；用 A 表示 \(\mathbf{M}_{n}(k)\) 的由上三角矩阵组成的子环，用 M 表示右 A-模 \(k^{n}\)。证明 M 是投射的且有限生成的，但 M 的位似环在其双交换子中不稠密（参见 VIII, 第 88 页，定理 3 及 VIII, 第 88 页，注）。特别地，M 不是平衡的。
\item
  Z-模 Q 的双交换子是什么？证明 Q 是其双交换子上的单模，但 Z-模 Q 的位似环在其双交换子中不稠密。
\end{enumerate}

¶ 15) 设 M 是主理想整环上的一个挠模。证明 M 的位似环在其双交换子中稠密。

（化归为 M 是 A 的不可约元素 \(\pi\) 的 \(\pi\)-准素模的情形。然后证明 M 的每个有限子集 S 都含于一个直因子子模中，该子模是一个有限生成模与有限个不可分解可除模的直和。借助 VII, §2, 第 54 页习题 3 及 VII, §2, 第 56 页习题 8，对 S 的基数用归纳法。）

\begin{enumerate}
\def\labelenumi{\arabic{enumi})}
\setcounter{enumi}{15}
\tightlist
\item
  设 D 是一个体，V 是 D 上一个维数 \(\geqslant 2\) 的向量空间，A 是 \(\operatorname{End}_{\mathrm{D}}(\mathrm{V})\) 的一个子环。
\end{enumerate}

\begin{enumerate}
\def\labelenumi{\alph{enumi})}
\item
  假设对每个由四个元素组成的组 \((x, x', y, y')\)（其中 x 与 \(x'\) 线性无关），存在 A 的元素 u，使得 \(u(x) = y\) 且 \(u(x') = y'\)。证明 A-模 V 的交换子是 D（与 V 的位似环等同），且这个模的位似环在其双交换子中稠密。
\item
  给出环 \(A \subset \operatorname{End}_{\mathrm{D}}(\mathrm{V})\) 的一个例子，使 V 是单 A-模但其交换子异于 D（考虑一个含 D 的体 E、E 上的一个向量空间 V 以及环 \(A = \operatorname{End}_{\mathrm{E}}(\mathrm{V})\)）。
\end{enumerate}

\begin{enumerate}
\def\labelenumi{\arabic{enumi})}
\setcounter{enumi}{16}
\tightlist
\item
  设 K 是一个代数闭交换域（域），V 是 K 上一个有限维向量空间，G 是 \(\operatorname{Aut}(V)\) 的一个子幺半群，使得 K{[}G{]}-模 V 是单模。证明若标量 \(\operatorname{Tr}(g)\)（\(g \in G\)）组成的集合有限，则 G 有限（由伯恩赛德定理推出 G 含一个基 \((g_{i})_{i \in I}\)（\(\operatorname{End}(V)\) 的），并考虑映射 \(g \mapsto (\operatorname{Tr}(gg_{i}))\)（从 G 到 \(K^{I}\)））。
\end{enumerate}

¶ 18) 若群 G 的每个具有有限生成族的子群都有限，则称 G 为局部有限的。

\begin{enumerate}
\def\labelenumi{\alph{enumi})}
\item
  设 G 是一个群，H 是 G 的一个正规子群。证明若 H 与 G/H 都局部有限，则 G 亦然。
\item
  由 a) 推出，一个其元素均为有限阶的可解群是局部有限的。
\item
  设 K 是一个交换域（域），V 是 K 上一个有限维向量空间，G 是 Aut(V) 的一个子群，它有一个由有限阶元素组成的有限生成族。证明 G 的元素的阶是有界的。
\end{enumerate}

（通过考虑由 G 的一个有限生成族的系数生成的子域 K，化归为 K 是素域 P 的纯超越扩张的情形。然后证明：对 \(\mathbf{GL}_{n}(\mathrm{K})\) 的每个有限阶元素 s，我们有 \(\varphi(s) \leqslant n\)（当 P = Q），且 \(s \leqslant p^{n} - 1\)（当 \(P = F_{p}\)）。）

\begin{enumerate}
\def\labelenumi{\alph{enumi})}
\setcounter{enumi}{3}
\tightlist
\item
  证明 Aut(V) 的每个其元素均为有限阶的子群 H 都是局部有限的。\(^{(1)}\)
\end{enumerate}

（化归为 K 代数闭的情形，并对 \(\dim(V)\) 用归纳法。设 G 是 H 的一个具有有限生成族的子群。若 V 是单 K{[}G{]}-模，应用习题 17；若 V 含一个在 G 下稳定的非平凡子空间 W，考虑 G 在 \(\operatorname{Aut}(V) \times \operatorname{Aut}(V/W)\) 中的像，并应用 a) 与 b)。）

¶ 19) 设 A 是一个环，使得 A-模 \(A_{s}\) 是一族有限的左理想 \((\mathfrak{I}_{i})_{1\leqslant i\leqslant n}\)（作为 A-模彼此同构）的直和。

\begin{enumerate}
\def\labelenumi{\alph{enumi})}
\item
  证明在 A 中存在一族元素 \(e_{i,j}\)（\(1 \leqslant i \leqslant n\)，\(1 \leqslant j \leqslant n\)），使得 \(e_{i,j}e_{h,k} = \delta_{j,h}e_{i,k}\)，其中 \(\delta_{j,h}\) 表示克罗内克符号，且 \(I_i = Ae_{i,i}\)（参见 VIII, 第 39 页，习题 5 及 I, §8, 第 173 页，习题 12）。反之，若在 A 中存在一族 \(n^2\) 个元素 \(e_{i,j}\)，使得 \(e_{i,j}e_{h,k} = \delta_{j,h}e_{i,k}\) 且 \(1 = \sum_{i=1}^{n} e_{i,i}\)，则诸左理想 \(Ae_{i,i}\) 作为 A-模同构（参见 VIII, 第 39 页，习题 5）。此外，若 B 是 A 中与所有 \(e_{i,j}\) 交换的元素组成的子环，则 A 同构于矩阵环 \(\mathbf{M}_n(\mathbf{B})\)，且 B 同构于诸 A-模 \(Ae_{i,i}\) 中每一个的交换子之反环（利用 VIII, 第 78 页及 VIII, 第 65 页，推论）。
\item
  设 M 是一个 A-模。证明诸 \(N_{i} = e_{i,i}M\) 是 B-模，且 M 作为 B-模是诸 \(N_{i}\) 的直和；此外，诸 B-模 \(N_{i}\) 两两同构（考虑映射 \(x \mapsto e_{i,1}x\) 与 \(y \mapsto e_{1,i}y\)），且诸 \(N_{i}\) 的零化子等于 B 与 M 在 A 中的零化子之交。反之，对每个 B-模 N，在 n 个同构于 N 的 B-模的直和 M 上定义一个 A-模结构，使 \(e_{1,1}M\) 同构于 N。
\item
  对 \(N_{1}\) 的每个 B-子模 P，指定 M 的 A-子模 \(\sum_{i=1}^{n} e_{i,1}P\)。证明这定义了从 \(N_{1}\) 的 B-子模组成的集合到 M 的 A-子模组成的集合的双射，且它对包含关系严格递增。
\end{enumerate}

